annihilated by @@M0@@, hence of @@M1@@-Modules, and we have:
@@M2@@
giving the result. We thus define a function on the category of coherent @@M3@@-Modules with values in @@M4@@ by assigning to every coherent @@M5@@-Module the class of one of its extensions to @@M6@@.

This function is additive. Indeed, let:
@@M7@@
be an exact sequence of coherent @@M8@@-Modules. Let @@M9@@ be an extension of @@M10@@ to @@M11@@; then there exists a sub-Module @@M12@@ of @@M13@@ extending @@M14@@, and @@M15@@ extends @@M16@@. It therefore factors through @@M17@@, defining the required homomorphism @@M18@@. One checks immediately that it is inverse to @@M19@@.

\sgaStableTarget{sga6.e09.prop.1-2}\textbf{Proposition 1.2.} \uline{Let} @@M20@@ \uline{be two noetherian schemes}, \uline{and} @@M21@@ \uline{a} flat \uline{morphism from} @@M22@@ \uline{to} @@M23@@. \uline{Suppose given a family of elements} @@M24@@ \uline{of} @@M25@@, \uline{and consider} \uline{the homomorphism}:
@@M26@@
\uline{defined by}:
@@M27@@
\uline{where} @@M28@@ \uline{is the} ``\uline{inverse-image}'' \uline{homomorphism from} @@M29@@ \uline{to} @@M30@@, \uline{defined by} IV 2.12, \uline{and where} @@M31@@ \uline{is regarded as a} @@M32@@\uline{-module by} IV 2.10. \uline{Suppose that for every point} @@M33@@, \uline{the images of the} @@M34@@ \uline{in} @@M35@@ \uline{under} \uline{the homomorphism}:
@@M36@@
\uline{generate} @@M37@@ \uline{as an abelian group}. \uline{Then} @@M38@@ \uline{is surjective}.
